\section{Теоретические основы визуальной навигации}

\subsection{Введение в ArUco и AprilTag}
Визуальная навигация — это способ определения положения робота в пространстве с помощью анализа изображений. В отличие от GPS, который работает только на улице и имеет погрешность в несколько метров, визуальные маркеры позволяют дрону позиционироваться с точностью до сантиметра даже в закрытых помещениях.

\textbf{ArUco} (Augmented Reality University of Cordoba) и \textbf{AprilTag} — это самые популярные библиотеки фидуциальных (опорных) маркеров. Они представляют собой черно-белые квадраты, которые содержат закодированную информацию (ID маркера) и позволяют алгоритмам компьютерного зрения однозначно определить ориентацию маркера в пространстве.

\subsection{Принцип работы алгоритма обнаружения}
Процесс обнаружения маркера на изображении состоит из нескольких этапов. Понимание этого конвейера (pipeline) важно для диагностики проблем, когда дрон <<не видит>> маркеры.

\subsubsection{1. Предварительная обработка (Preprocessing)}
Камера передает цветное изображение (RGB). Алгоритм сначала переводит его в оттенки серого (Grayscale), так как цвет не несет информации для черно-белых маркеров.

Затем применяется \textbf{адаптивная бинаризация} (Adaptive Thresholding). В отличие от простого порога (где все пиксели ярче 127 становятся белыми, а темнее — черными), адаптивный метод вычисляет порог для каждой небольшой области изображения. Это позволяет находить маркеры даже при неравномерном освещении (например, когда одна часть маркера в тени).

\begin{codefiletitle}[python]{Пример бинаризации (OpenCV)}{python/examples/advanced_navigation/binarization_example.py}
\end{codefiletitle}

\subsubsection{2. Поиск контуров (Contour Extraction)}
Алгоритм ищет границы связных областей белых пикселей на черном фоне. Из всех найденных контуров отбираются только те, которые:
\begin{itemize}
    \item Являются замкнутыми.
    \item Аппроксимируются четырехугольником (имеют 4 вершины).
    \item Являются выпуклыми (Convex).
    \item Имеют достаточный размер (площадь).
\end{itemize}

\subsubsection{3. Исправление перспективы (Perspective Warping)}
Найденный четырехугольник может быть искажен перспективой (выглядеть как трапеция или ромб), если камера смотрит на него под углом. Алгоритм вычисляет матрицу гомографии, чтобы <<выпрямить>> изображение и получить идеальный квадрат (Canonical Marker Image).

\subsubsection{4. Декодирование (Decoding)}
Выпрямленный квадрат разбивается на сетку (например, $6 \times 6$ клеток). Каждая клетка анализируется: черная — 0, белая — 1. Полученная матрица битов сверяется со словарем (Dictionary).

\subsection{Коррекция ошибок и Расстояние Хэмминга}
Маркеры ArUco спроектированы так, чтобы быть устойчивыми к ошибкам считывания. Даже если несколько пикселей (битов) будут распознаны неверно (из-за блика или грязи), алгоритм все равно сможет узнать маркер.

Это достигается за счет избыточности информации. Каждый маркер имеет параметр \textbf{Hamming Distance} (Расстояние Хэмминга) — минимальное количество битов, на которое отличаются любые два маркера в словаре.
\begin{itemize}
    \item Если словарь имеет Hamming Distance = 4, система может обнаружить до 3 ошибок и исправить до 1 ошибки.
    \item Чем больше размер словаря (например, $6 \times 6$ против $4 \times 4$), тем больше информации можно закодировать и тем выше устойчивость к ложным срабатываниям.
\end{itemize}

\subsection{Сравнение типов словарей (Dictionaries)}
При выборе словаря нужно искать баланс между количеством доступных ID и размером маркера (читаемостью).

\begin{table}[H]
\centering
\begin{tabularx}{\textwidth}{|l|X|X|X|}
\hline
\textbf{Словарь} & \textbf{Размер сетки} & \textbf{Кол-во ID} & \textbf{Применение} \\ \hline
DICT\_4X4\_50 & $4 \times 4$ бита & 50 & Дальние дистанции, низкое разрешение. Меньше деталей — легче распознать. \\ \hline
DICT\_5X5\_1000 & $5 \times 5$ бит & 1000 & Стандарт для большинства задач. \\ \hline
DICT\_6X6\_250 & $6 \times 6$ бит & 250 & Высокая надежность, используется в GeoSkan Pioneer. \\ \hline
DICT\_APRILTAG\_36h11 & $6 \times 6$ бит & 587 & Промышленный стандарт, очень высокая устойчивость к ложным срабатываниям. \\ \hline
\end{tabularx}
\caption{Характеристики словарей маркеров}
\end{table}

\subsection{Математика позиционирования (PnP Problem)}
После того как маркер найден и идентифицирован, нужно определить положение камеры относительно него. Эта задача называется \textbf{Perspective-n-Point (PnP)}.

Мы знаем:
\begin{enumerate}
    \item Координаты 4 углов маркера в его собственной системе координат (3D object points): $(0,0,0), (S,0,0), (S,S,0), (0,S,0)$, где $S$ — размер маркера.
    \item Координаты этих же 4 углов на изображении камеры (2D image points) в пикселях.
    \item Внутренние параметры камеры (Camera Matrix и Distortion Coefficients), полученные при калибровке.
\end{enumerate}

Функция \texttt{cv2.solvePnP} решает систему уравнений и возвращает два вектора:
\begin{itemize}
    \item \texttt{rvec} (Rotation Vector) — вектор вращения (в формате Родригеса), показывающий наклон камеры.
    \item \texttt{tvec} (Translation Vector) — вектор смещения $(x, y, z)$, показывающий позицию маркера относительно центра камеры.
\end{itemize}

\begin{alertbox}{Важно!}
Вектор \texttt{tvec} показывает координаты \textit{маркера в системе координат камеры}. Чтобы получить координаты \textit{дрона в системе координат маркера} (для навигации), нужно инвертировать эту трансформацию:
$$ R = Rodrigues(rvec) $$
$$ Position_{drone} = -R^T \cdot tvec $$
\end{alertbox}

\subsection{Практические задания}

\begin{taskbox}[Задание 1: Анализ словарей]
Сгенерируйте маркер ID=1 из словарей 4x4, 5x5 и 6x6 одного физического размера (например, 200 пикселей).
1. Распечатайте их.
2. Отойдите с камерой на максимальное расстояние.
3. Какой маркер перестанет распознаваться первым? Почему?
\end{taskbox}

\begin{taskbox}[Задание 2: Влияние освещения]
Запустите скрипт обнаружения маркера.
1. Постепенно затемняйте комнату. На каком уровне освещенности маркер пропадет?
2. Посветите фонариком прямо на маркер (создайте блик). Как это влияет на бинаризацию (посмотрите окно с ч/б изображением)?
\end{taskbox}

\begin{questionbox}[Контрольные вопросы]
\begin{enumerate}
    \item Почему для навигации не используются QR-коды? (Подсказка: сравните сложность детектирования угла наклона).
    \item Что такое <<адаптивная бинаризация>> и зачем она нужна?
    \item Если дрон видит маркер под углом 45 градусов, как алгоритм понимает, что это квадрат, а не ромб?
    \item В чем физический смысл вектора \texttt{tvec}?
\end{enumerate}
\end{questionbox}
